\documentclass[12pt,a4paper]{article}

\usepackage[utf8]{inputenc}
\usepackage[T2A]{fontenc}
\usepackage[ukrainian]{babel}
\usepackage{geometry}
\usepackage{setspace}
\usepackage{array}
\usepackage{longtable}
\usepackage{booktabs}
\usepackage{ragged2e}
\usepackage{indentfirst}
\usepackage{titlesec}
\usepackage{tocloft}

\geometry{
    left=30mm,
    right=15mm,
    top=20mm,
    bottom=20mm
}

\usepackage[
  unicode,
  pdftitle={Технічне завдання},
  pdfauthor={Давидчук Артем}
]{hyperref}

\setlength{\parindent}{1.25cm}
\onehalfspacing

\titleformat{\section}
    {\normalfont\bfseries}
    {\thesection.}
    {0.5em}
    {}

\renewcommand{\contentsname}{Зміст}

\begin{document}

    \begin{center} \textbf{\Large Технічне завдання на курсову роботу} \end{center}

    \vspace{0.5cm}

    \tableofcontents

    \newpage

    \section{НАЙМЕНУВАННЯ ТА ОБЛАСТЬ ВИКОРИСТАННЯ}

    Це технічне завдання поширюється на розробку програмного додатка
    \textbf{``Система діалогової взаємодії з великими мовними моделями з веб- та консольним інтерфейсами''}.

    Область використання розроблюваного програмного продукту:
    організація діалогової взаємодії користувача з мовною моделлю,
    ведення історії чатів, запуск застосунку в різних режимах роботи,
    а також використання в навчальних, демонстраційних та дослідницьких цілях
    у галузі інформаційних технологій, штучного інтелекту та програмної інженерії.

    Програмний додаток призначений для локального або мережевого використання
    на персональному комп'ютері та має забезпечувати зручну взаємодію з системою
    як через веб-інтерфейс користувача у браузері, так і через інтерфейс командного рядка.

    \section{ПРИЧИНИ ДЛЯ РОЗРОБКИ}

    Підставою для розробки є завдання на виконання курсової роботи
    за освітньо-професійною програмою
    \textbf{<<Комп'ютерні системи та мережі>>}
    спеціальності

    \noindent
    \textbf{123 <<Комп'ютерна інженерія>>}.

    Необхідність розробки зумовлена зростанням актуальності систем,
    що забезпечують взаємодію людини з інтелектуальними програмними засобами,
    зокрема з великими мовними моделями, а також потребою у створенні
    зручного, гнучкого та доступного програмного інструменту,
    який підтримує декілька режимів роботи -- веб-інтерфейс та консольний інтерфейс.

    \section{МЕТА ТА ПРИЗНАЧЕННЯ РОЗРОБКИ}

    Метою даного проєкту є розробка програмного додатка для діалогової взаємодії
    з великою мовною моделлю, який забезпечує:

    \begin{itemize}
        \item запуск у двох режимах роботи: через веб-інтерфейс користувача у браузері та через командний інтерфейс (CLI);
        \item створення нового чату та вибір наявного чату;
        \item введення, передавання та відображення повідомлень у межах діалогу;
        \item збереження історії листування;
        \item зручність використання та зрозумілу організацію інтерфейсу.
    \end{itemize}

    Призначенням розробки є створення універсального програмного засобу,
    який може бути використаний як демонстраційна, навчальна або базова прикладна система
    для роботи з мовними моделями та подальшого розширення функціональних можливостей.

    \section{ДЖЕРЕЛА РОЗРОБКИ}

    Джерелами розробки є:

    \begin{itemize}
        \item науково-технічна література з програмної інженерії, штучного інтелекту та побудови діалогових систем;
        \item довідкова документація з мови програмування Python;
        \item документація до бібліотек, що використовуються для реалізації веб- та консольного інтерфейсів;
        \item публікації в мережі Інтернет, присвячені побудові LLM-систем, клієнтських застосунків та інтерфейсів користувача;
        \item матеріали з проєктування архітектури програмного забезпечення та організації збереження даних.
    \end{itemize}

    \section{ТЕХНІЧНІ ВИМОГИ}

    \subsection{Вимоги до розробленого продукту}

    Розроблюваний програмний додаток повинен відповідати таким вимогам:

    \begin{enumerate}
        \item Програма повинна запускатися з головного файла \texttt{app.py}.
        \item Має бути реалізовано два режими роботи:
        \begin{itemize}
            \item стандартний запуск -- запуск веб-інтерфейсу користувача з доступом через браузер;
            \item запуск з параметром \texttt{-t} -- перехід до режиму командного рядка.
        \end{itemize}
        \item У режимі веб-інтерфейсу має бути забезпечено:
        \begin{itemize}
            \item відображення списку чатів;
            \item створення нового чату;
            \item відкриття існуючого чату;
            \item відображення історії повідомлень;
            \item введення нового повідомлення та отримання відповіді.
        \end{itemize}
        \item У консольному режимі має бути забезпечено:
        \begin{itemize}
            \item створення нового чату;
            \item вибір існуючого чату;
            \item обмін повідомленнями у текстовому режимі;
            \item коректне завершення сеансу роботи.
        \end{itemize}
        \item Повинно бути реалізоване збереження історії чатів у файлах або локальному сховищі даних.
        \item Інтерфейс програми повинен бути логічно структурованим, зрозумілим та зручним для користувача.
        \item Програма повинна коректно обробляти помилки введення, збої під час запуску та некоректні параметри командного рядка.
        \item Архітектура програмного забезпечення повинна передбачати можливість подальшого розширення функціональності.
    \end{enumerate}

    \subsection{Вимоги до програмного забезпечення}

    Для функціонування програмного додатка необхідні:

    \begin{itemize}
        \item операційна система: Windows 10/11, Linux або інша сучасна операційна система з підтримкою Python;
        \item інтерпретатор мови програмування Python версії 3.10 або вище;
        \item сучасний веб-браузер з підтримкою HTML5, CSS та JavaScript;
        \item бібліотеки та модулі, необхідні для реалізації веб-інтерфейсу користувача, логіки програми та роботи із збереженням чатів;
        \item засоби розробки та налагодження програмного забезпечення (за потреби).
    \end{itemize}

    \subsection{Вимоги до апаратної частини}

    Для роботи програмного додатка необхідний персональний комп'ютер
    із такими мінімальними характеристиками:

    \begin{itemize}
        \item процесор архітектури x86\_64;
        \item оперативна пам'ять -- не менше 8 ГБ;
        \item вільне місце на носії даних -- не менше 1 ГБ;
        \item монітор, клавіатура та пристрій введення;
        \item за потреби -- підключення до мережі Інтернет для роботи із зовнішніми сервісами або моделями.
    \end{itemize}

    \section{ЕТАПИ РОЗРОБКИ}

    \begin{longtable}{|>{\centering\arraybackslash}m{1.5cm}|>{\RaggedRight\arraybackslash}m{10.5cm}|>{\centering\arraybackslash}m{3cm}|}
        \hline
        \textbf{№} & \textbf{Етап розробки} & \textbf{Дата} \\
        \hline
        6.1 & Отримання теми та завдання & 15.02.2026\\
        \hline
        6.2 & Підбір та вивчення літератури за тематикою діалогових систем, LLM та веб-інтерфейсів користувача & 15.03.2026\\
        \hline
        6.3 & Формування технічного завдання & 15.03.2026\\
        \hline
        6.4 & Аналіз предметної області та огляд існуючих підходів до побудови діалогових застосунків & \\
        \hline
        6.5 & Проєктування структури програмного додатка, режимів веб-інтерфейсу та CLI, а також механізму збереження чатів & \\
        \hline
        6.6 & Розробка програмного додатка & \\
        \hline
        6.7 & Тестування програмного додатка & \\
        \hline
        6.8 & Оформлення пояснювальної записки & \\
        \hline
        6.9 & Подання курсового проєкту (роботи) на перевірку & \\
        \hline
        6.10 & Захист курсового проєкту (роботи) & \\
        \hline
    \end{longtable}

\end{document}